\exercisesection{}

Unless stated otherwise, the Banach spaces and Hilbert spaces below are assumed to be complex.

\begin{enumerate}
\def\labelenumi{\arabic{enumi})}
\tightlist
\item
  \begin{enumerate}
  \def\labelenumii{\alph{enumii})}
  \tightlist
  \item
    There exist vector spaces E, F, G and partial operators u, v from E to F, and w from F to G, such that \(w \circ (u + v) \neq w \circ u + w \circ v\) .
  \end{enumerate}
\end{enumerate}

\begin{enumerate}
\def\labelenumi{\alph{enumi})}
\setcounter{enumi}{1}
\item
  Give an example of a closable operator u on a topological vector space such that the domain of the closure of u is strictly contained in the closure of the domain of u.
\item
  Give an example of closed operators u and v on a Hilbert space such that u is continuous but \(u \circ v\) is not closed.
\end{enumerate}

\begin{enumerate}
\def\labelenumi{\arabic{enumi})}
\setcounter{enumi}{1}
\tightlist
\item
  Let E and F be Banach spaces, and let G be a subspace of E. Let u be a finite-rank linear map from G into F. Suppose that u is not continuous on G.
\end{enumerate}

\begin{enumerate}
\def\labelenumi{\alph{enumi})}
\item
  The partial operator from E to F with domain G defined by u is not closable.
\item
  Suppose that E and G are Hilbert spaces. Compute the adjoint of the partial operator defined by u.
\end{enumerate}

\begin{enumerate}
\def\labelenumi{\arabic{enumi})}
\setcounter{enumi}{2}
\tightlist
\item
  Let K = R or C. Let E, F be topological vector spaces over K and u an operator with dense domain from E into F. Let D' be the set of continuous linear forms \(\lambda \in F'\) such that there exists a continuous linear form \(\mu \in E'\) coinciding with \(\lambda \circ u\) on the domain of u.
\end{enumerate}

\begin{enumerate}
\def\labelenumi{\alph{enumi})}
\tightlist
\item
  The set D' is a vector subspace of F', the linear form \(\mu\) is determined in a unique way by \(\lambda\) , and the map \(\lambda \mapsto \mu\) from D' to E' is linear.
\end{enumerate}

One calls the transpose of u, and denotes \(^{t}u\) the partial operator from F' into E' whose domain is D' and which is given by \(\lambda\mapsto\mu\) for \(\lambda\in D'\) .

\begin{enumerate}
\def\labelenumi{\alph{enumi})}
\setcounter{enumi}{1}
\item
  Suppose that E and F are locally convex. The operator u is closable if and only if the duality \(\mathrm{F} \times \mathrm{dom}(^{t}u) \to \mathrm{K}\) is separating.
\item
  Suppose that E and F are locally convex and that F is semi-reflexive. The operator u is closable if and only if \(^{t}u\) has dense domain.
\end{enumerate}

\begin{enumerate}
\def\labelenumi{\arabic{enumi})}
\setcounter{enumi}{3}
\item
  Let E be a complex Hilbert space and u an injective self-adjoint partial operator on E. The partial operator \(u^{-1}\) is self-adjoint.
\item
  Let E be a complex Banach space. A continuous semigroup on E is called a map \(s\colon\mathbf{R}_{+}\to\mathcal{L}(\mathrm{E})\) such that
\end{enumerate}

\begin{enumerate}
\def\labelenumi{(\roman{enumi})}
\item
  \(s(0) = 1_{\mathrm{E}}\) ;
\item
  The map from \(R_{+} \times E \to E\) defined by \((t, x) \mapsto s(t)x\) is continuous;
\item
  For all \(t_1\) and \(t_2\) in \(\mathbf{R}_{+}\) , we have \(s(t_1 + t_2) = s(t_1)s(t_2)\) .
\end{enumerate}

We consider a continuous semigroup s on E.

\begin{enumerate}
\def\labelenumi{\alph{enumi})}
\tightlist
\item
  There exist real numbers C ≥ 0 and M ≥ 0 such that
\end{enumerate}

\[
\| s (t) \| \leqslant \mathrm{Ce} ^ {\mathrm{Mt}}
\]

for all \(t \in R_{+}\) . (Let a \textgreater{} 0; show that \(\|s(t)\|\) is bounded for \(t \in [0, a]\) .) We fix such numbers in the rest of the exercise.

\begin{enumerate}
\def\labelenumi{\alph{enumi})}
\setcounter{enumi}{1}
\item
  Let D be the set of \(x \in E\) such that the mapping \(\psi_{x}\) of \(R_{+}\) to E defined by \(t \mapsto s(t)x\) is differentiable at 0. The space D is dense in E; the map u from D into E defined by \(u(x) = -\psi_{x}'(0)\) is a closed partial operator on E, which is called the infinitesimal generator of the semigroup s.
\item
  Let \(\lambda\) a complex number such that \(\mathcal{R}(\lambda) < -\mathrm{M}\) . The formula
\end{enumerate}

\[
v _ {\lambda} = - \int_ {\mathbf {R} _ {+}} e ^ {\lambda t} s (t) d t
\]

defines an endomorphism of E.

\(d)\) The image of \(v_{\lambda}\) is contained in \(\operatorname{dom}(u)\) and we have \(u \circ v_{\lambda} = -1_{\mathrm{E}} - \lambda v_{\lambda}\) .

\begin{enumerate}
\def\labelenumi{\alph{enumi})}
\setcounter{enumi}{4}
\tightlist
\item
  One has \(u \circ v_{\lambda} = v_{\lambda} \circ u\) .
\end{enumerate}

\(f)\) The complex number \(\lambda\) belongs to the resolvent set of \(u\) and we have \(v_{\lambda} = -\mathrm{R}(u,\lambda)\) .

\begin{enumerate}
\def\labelenumi{\alph{enumi})}
\setcounter{enumi}{6}
\tightlist
\item
  One has
\end{enumerate}

\[
\| v _ {\lambda} ^ {n} \| \leqslant \frac {\mathrm{C}}{| \mathcal {R} (\lambda) + \mathrm{M} | ^ {n}}
\]

for all \(n \in N\) .

\begin{enumerate}
\def\labelenumi{\arabic{enumi})}
\setcounter{enumi}{5}
\tightlist
\item
  Let E be a complex Banach space and u a closed partial operator with dense domain on E. Suppose there exist real numbers \(C \geqslant 0\) and \(M \geqslant 0\) such that
\end{enumerate}

\begin{enumerate}
\def\labelenumi{(\roman{enumi})}
\item
  The set {]}−∞, −M{[} is contained in the resolvent set of u;
\item
  For every integer \(n \in \mathbf{N}\) and every real number \(\lambda < -\mathrm{M}\) , we have
\end{enumerate}

\[
\left\| \mathrm{R} (u, \lambda) ^ {n} \right\| \leqslant \mathrm{C} | \lambda + \mathrm{M} | ^ {- n}.
\]

\begin{enumerate}
\def\labelenumi{\alph{enumi})}
\tightlist
\item
  For every real number \(y < -\mathrm{M}\) , we denote
\end{enumerate}

\[
u _ {y} = - y 1 _ {\mathrm{E}} + y ^ {2} \mathrm{R} (u, y) \in \mathcal {L} (\mathrm{E}).
\]

The endomorphisms \(u_{y}\) commute pairwise; for every \(x \in \text{dom}(u)\) , we have \(u_{y}(x) \to u(x)\) as \(y \to -\infty\) .

\begin{enumerate}
\def\labelenumi{\alph{enumi})}
\setcounter{enumi}{1}
\item
  The map from \(R_{+}\) in \(\mathcal{L}(E)\) defined by \(s_{y}(t)=e^{-tu_{y}}\) is a continuous semigroup on E such that \(\|s_{y}(t)\|\leqslant Ce^{Mt}\) for all \(t\in R_{+}\) .
\item
  For all \(y_{1}, y_{2}\) and \(t\) in \(\mathbf{R}_{+}^{*}\) , and every \(x\) in \(\mathrm{dom}(u)\) , we have
\end{enumerate}

\[
\left\| s _ {y _ {1}} (t) x - s _ {y _ {2}} (t) x \right\| \leqslant t \mathrm{C} ^ {2} \left\| u _ {y _ {1}} (x) - u _ {y _ {2}} (x) \right\|.
\]

\(d)\) Let \(t \in \mathbf{R}_{+}\) . For every \(x\) in E, the function defined on \(\mathbf{R}_{+}^{*}\) by \(y \mapsto s_{y}(t)x\) converges as \(y \to -\infty\) to an element \(s(t)x\) of E. The map \(s(t)\) is linear and continuous with norm \(\leqslant \mathrm{Ce}^{\mathrm{M}t}\) .

\begin{enumerate}
\def\labelenumi{\alph{enumi})}
\setcounter{enumi}{4}
\item
  The map \(t \mapsto s(t)\) is a continuous semigroup on E.
\item
  The partial operator u is the infinitesimal generator of the continuous semigroup s (“Hille–Yosida theorem”). (First verify that the infinitesimal generator of s is an extension of u.)
\end{enumerate}

\(g)\) Let \(\widetilde{s}\) a continuous semigroup whose \(u\) is the infinitesimal generator. We have \(\widetilde{s} = s\) .

\begin{enumerate}
\def\labelenumi{\arabic{enumi})}
\setcounter{enumi}{6}
\tightlist
\item
  Let E be a complex Banach space and u a partial operator on E. We say that u is accretive if, for every \(x \in \text{dom}(u)\) there exists a continuous linear form \(\ell \in E'\) such that
\end{enumerate}

\begin{enumerate}
\def\labelenumi{(\roman{enumi})}
\item
  \(\mathcal{R}(\ell(u(x))) \geqslant 0;\)
\item
  One has \(\|\ell\|=\|x\|\) and \(\ell(x)=\|x\|^{2}\) .
\end{enumerate}

We say that a continuous semigroup \(s\) on E is contractive if \(\|s(t)\|\leqslant1\) for all \(t\in\mathbf{R}_{+}\) .

\begin{enumerate}
\def\labelenumi{\alph{enumi})}
\item
  The partial operator \(u\) is the infinitesimal generator of a continuous contractive semigroup on E if and only if \(u\) is accretive and if there exists \(\lambda > 0\) such that the partial operator \(\lambda 1_{\mathrm{E}} + u\) is surjective. (To show that the condition is necessary, let \(s\) be the contractive semigroup whose infinitesimal generator is \(u\) ; for \(x \in \operatorname{dom}(u)\) , consider the function defined by \(t \mapsto \ell(s(t)x)\) ; it is differentiable at 0, and its derivative at 0 is \(-\mathcal{R}(\ell(u(x)))\) .)
\item
  Suppose that u is closed and that \(u\) and \(^{t}u\) are accretive. Then u is the infinitesimal generator of a continuous contractive semigroup on E.
\end{enumerate}

\begin{enumerate}
\def\labelenumi{\arabic{enumi})}
\setcounter{enumi}{7}
\tightlist
\item
  Let E be a complex Banach space. Let s be a continuous contractive semigroup on E, and let u be its infinitesimal generator. We have
\end{enumerate}

\[
\lim _ {n \to + \infty} \left(1 + \frac {t}{n} u\right) ^ {- n} x = s (t) x
\]

for all \(x \in E\) .

\begin{enumerate}
\def\labelenumi{\arabic{enumi})}
\setcounter{enumi}{8}
\tightlist
\item
  \begin{enumerate}
  \def\labelenumii{\alph{enumii})}
  \tightlist
  \item
    Find Hilbert spaces E and F and partial operators u and v from E into F such that u and \(u + v\) are densely defined and \(u^{*} + v^{*}\) is different from \((u + v)^{*}\) .
  \end{enumerate}
\end{enumerate}

\begin{enumerate}
\def\labelenumi{\alph{enumi})}
\setcounter{enumi}{1}
\item
  Find Hilbert spaces E, F, and G and densely defined partial operators u and v from E into F and from F into G, respectively, such that \(v \circ u\) is densely defined and \(u^{*} \circ v^{*}\) is different from \((v \circ u)^{*}\) .
\item
  Find a Hilbert space E and partial self-adjoint operators u and v on E such that \(u + v\) is densely defined but is not self-adjoint.
\end{enumerate}

\begin{enumerate}
\def\labelenumi{\arabic{enumi})}
\setcounter{enumi}{9}
\item
  Give an example of multiplication operators. \(m_{g_{1}}\) and \(m_{g_{2}}\) such that \(m_{g_{1}} \circ m_{g_{2}} \neq m_{g_{1}g_{2}}\) .
\item
  We identify the space \(\mathcal{D}(]0,1[)\) with a subspace of \(\mathrm{L}^2 ([0,1])\) Let \(u\) the partial operator on \(\mathrm{L}^2 ([0,1])\) with domain \(\mathcal{D}(]0,1[)\) such that \(u(f) = i^{-1}f'\) for every function \(f \in \mathcal{D}(]0,1[)\) .
\end{enumerate}

\begin{enumerate}
\def\labelenumi{\alph{enumi})}
\item
  The partial operator \(u\) is closable and \(\overline{u}\) is symmetric.
\item
  The eigenspaces of \(\overline{u}\) relative to i and to -i are not orthogonal.
\end{enumerate}

\begin{enumerate}
\def\labelenumi{\arabic{enumi})}
\setcounter{enumi}{11}
\item
  Let \(S \subset \mathbf{C}\) a closed subset. Give an example of a closed operator with dense domain \(u\) on a Hilbert space \(E\) such that \(\operatorname{Sp}(u) = S\) .
\item
  Let \(u\) a non-closable partial operator on a Hilbert space E. There does not exist a complex number \(\lambda\) such that the linear map from \(\text{dom}(u)\) to E defined by \(x \mapsto \lambda x - u(x)\) is bijective and has a continuous inverse.
\item
  Let \(u\) a symmetric partial operator on a complex Hilbert space E. The map \(\lambda \mapsto \dim \operatorname{Ker}(\lambda 1_{\mathrm{E}} - u^{*})\) , of \(\mathbf{C}\) in \(\overline{\mathbf{N}}\) is locally constant on the resolvent set of \(u\) .
\item
  Let E be the complex Hilbert space \(\mathbf{C}^2\) and let \(u \in \mathcal{L}(\mathrm{E})\) the endomorphism such that \(u(ae_1 + be_2) = be_1\) .
\end{enumerate}

\begin{enumerate}
\def\labelenumi{\alph{enumi})}
\tightlist
\item
  Show that the spectrum of \(u\) is reduced to 0 and that, for all \(\lambda \in \mathbf{C}^*\) , we have the formula
\end{enumerate}

\[
\| \mathrm{R} (u, \lambda) \| ^ {2} = \frac {2}{1 + 2 | \lambda | ^ {2} - \sqrt {1 + 4 | \lambda | ^ {2}}}.
\]

Deduce the form of the pseudo-spectrum from this. \(\mathrm{PSp}_{\varepsilon}(u)\) for all \(\varepsilon > 0\) .

\begin{enumerate}
\def\labelenumi{\arabic{enumi})}
\setcounter{enumi}{15}
\item
  Find an example showing that the conclusion of Prop. 21 of IV, p. 251, is not valid in general for bounded components of the spectrum.
\item
  \begin{enumerate}
  \def\labelenumii{\alph{enumii})}
  \tightlist
  \item
    Let U be a connected open subset of C and let V ⊂ U be an open set. Let E be a complex Banach space. Let \(f\colon\) U → E be a holomorphic map, \(z_0 \in\) U and M \textgreater{} 0 such that \(\|f(z)\| \leqslant\) M for all \(z \in\) V and such that \(\|f(z_0)\| <\) M. We then have \(\|f(z)\| <\) M for all \(z \in\) U.
  \end{enumerate}
\end{enumerate}

\begin{enumerate}
\def\labelenumi{\alph{enumi})}
\setcounter{enumi}{1}
\tightlist
\item
  Let E be a complex Banach space and \(u \in \mathcal{L}(\mathrm{E})\) Let \(U \subset C\) an open subset of the unbounded connected component of \(\mathbf{C}-\mathrm{Sp}(u)\) If M \textgreater{} 0 is such that \(\|R(u, \lambda)\| \leqslant M\) for all \(\lambda \in U\) , then \(\|R(u, \lambda)\| < M\) for all \(\lambda \in U\) . In particular, if \(\mathbf{C}-\mathrm{Sp}(u)\) is connected, then the norm of the resolvent of u cannot be constant on a nonempty open subset of \(\mathbf{C}-\mathrm{Sp}(u)\) .
\end{enumerate}

\begin{enumerate}
\def\labelenumi{\arabic{enumi})}
\setcounter{enumi}{17}
\tightlist
\item
  Let E be the complex vector space of bounded sequences \((x_{k})_{k\in\mathbf{Z}}\) of complex numbers. For every \(x=(x_{k})\) in E, we denote
\end{enumerate}

\[
\| x \| _ {\mathrm{E}} = | x _ {0} | + \sup _ {k \in \mathbf {Z} - \{0 \}} | x _ {k} |.
\]

\begin{enumerate}
\def\labelenumi{\alph{enumi})}
\item
  The map \(x \mapsto \|x\|_{\mathrm{E}}\) is a norm on E; the space E equipped with this norm is a Banach space.
\item
  Let M \textgreater{} 2 be a real number. Let \(\lambda = (\lambda_k)_{k \in \mathbf{Z}}\) the sequence defined by \(\lambda_0 = M^{-1}\) and \(\lambda_k = 1\) if \(k \neq 0\) . We denote by u the linear map from E to E such that \(u((x_k)_{k \in \mathbf{Z}}) = (\lambda_k x_{k+1})_{k \in \mathbf{Z}}\) . This is a continuous endomorphism of E.
\item
  The open disk with center 0 and radius 1/M is contained in the resolvent set of u.
\end{enumerate}

\(d)\) For every \(\lambda\in\mathbf{C}\) such that \(|\lambda|<\inf(M^{-1},\frac{1}{2}-M^{-1})\) , we have \(\|\mathrm{R}(u,\lambda)\|\leqslant\mathrm{M}\) . (Use formula (6) of IV, p. 244).

\begin{enumerate}
\def\labelenumi{\alph{enumi})}
\setcounter{enumi}{4}
\tightlist
\item
  Let \(e_{0} \in E\) the sequence \((x_{k})_{k \in \mathbf{N}}\) defined by \(x_{0} = 1\) and \(x_{k} = 0\) if \(k \neq 0\) Let \(\lambda \in C\) such that \(|\lambda| < \inf(M^{-1}, \frac{1}{2} - M^{-1})\) . We have \(\|R(u, \lambda)e_{0}\|_{E} = M\) , hence \(\|R(u, \lambda)\| = M\) .
\end{enumerate}

(This example, like that of the following exercise, is due to E. Shargorodsky, “On the level sets of the resolvent norm of a linear operator”, Bull. London Math. Soc. 40 (2008), 493–504.)

\begin{enumerate}
\def\labelenumi{\arabic{enumi})}
\setcounter{enumi}{18}
\tightlist
\item
  Let E be the Hilbert space \(\ell_2(\mathbf{N}^*)\) Let \((\alpha_n)_{n \in \mathbf{N}^*}\) a sequence of real numbers \(\geqslant 2\) such that \(\alpha_n\) tends to \(+\infty\) as \(n \to +\infty\) . Set \(\beta_n = 1 + \alpha_n^{-1}\) for all \(n \in \mathbf{N}^*\) .
\end{enumerate}

\begin{enumerate}
\def\labelenumi{\alph{enumi})}
\tightlist
\item
  Let \(\mathrm{F}\subset \mathrm{E}\) the space of sequences \((x_{k})_{k\in \mathbf{N}^{*}}\in \mathrm{E}\) such that the series
\end{enumerate}

\[
\sum_ {n \in \mathbf {N} ^ {*}} \alpha_ {n} ^ {2} | x _ {2 n} | ^ {2}
\]

converges. It is a dense subspace of E.

\begin{enumerate}
\def\labelenumi{\alph{enumi})}
\setcounter{enumi}{1}
\tightlist
\item
  There exists a closed partial operator. \(u\) on E with domain F such that for every \(x = (x_{k})_{k \in \mathbf{N}^{*}}\) in F, we have \(u(x) = (y_{k})\) where
\end{enumerate}

\[
\binom{y _ {2 k - 1}}{y _ {2 k}} = \left( \begin{array}{c c} 0 & \alpha_ {k} \\ \beta_ {k} & 0 \end{array} \right) \binom{x _ {2 k - 1}}{x _ {2 k}}
\]

for all \(k \in \mathbf{N}^*\) .

\begin{enumerate}
\def\labelenumi{\alph{enumi})}
\setcounter{enumi}{2}
\tightlist
\item
  Let \(u_{n}\) the endomorphism
\end{enumerate}

\[
x \mapsto \left( \begin{array}{c c} 0 & \alpha_ {n} \\ \beta_ {n} & 0 \end{array} \right) x
\]

of the Hilbertian space \(C^{2}\) equipped with the canonical inner product. For every \(\lambda \in C\) such that \(|\lambda| < 1\) , the resolvent of \(u_{n}\) at \(\lambda\) is defined and we have

\[
\lim _ {n \to + \infty} \| \mathrm{R} (u _ {n}, \lambda) \| = 1.
\]

\begin{enumerate}
\def\labelenumi{\alph{enumi})}
\setcounter{enumi}{3}
\tightlist
\item
  For every \(n \in N^{*}\) and every \(\lambda \in C\) such that \(|\lambda| < \frac{1}{2}\) , we have \(\|R(u_{n}, \lambda)\| < 1\) .
\item
  The open disk with center 0 and radius \(\frac{1}{2}\) is contained in the resolvent set of u, and we have \(\|R(u,\lambda)\|=1\) for all \(\lambda\in C\) such that \(|\lambda|<\frac{1}{2}\) .
\end{enumerate}

\begin{enumerate}
\def\labelenumi{\arabic{enumi})}
\setcounter{enumi}{19}
\tightlist
\item
  If \(a\) and \(b\) are real numbers with \(a < b\) , we denote \(\mathcal{C}^2([a,b])\) the vector space of continuous functions on \([a,b]\) whose restriction to \(]a,b[\) is of class \(\mathrm{C}^2\) . We denote \(\mu\) Lebesgue measure.
\end{enumerate}

\begin{enumerate}
\def\labelenumi{\alph{enumi})}
\tightlist
\item
  Let \(q_{1}\) and \(q_{2}\) of real-valued functions in \(\mathcal{C}^{2}([a,b])\) such that \(q_{1} \geqslant q_{2}\) . For \(i \in \{1,2\}\) , let \(f_{i} \in \mathcal{C}^{2}([a,b])\) a solution of the equation
\end{enumerate}

\[
f _ {i} ^ {\prime \prime} + q _ {i} f _ {i} = 0.
\]

Suppose that \(f_{2}(a) = f_{2}(b) = 0\) and that \(f_{2} \neq 0\) If the function \(f_{1}\) does not vanish in \(]a, b[\) , then \(q_{1} = q_{2}\) and there exists \(c \in \mathbf{C}\) such that \(f_{1} = cf_{2}\) . (Reduce to the case where \(f_1\) and \(f_2\) are strictly positive on \(]a, b[;\) consider the Wronskian \(w = f_1f_2' - f_1'f_2\) and prove that \(w(a) = w(b)\) then that \(w = 0.\) )

\begin{enumerate}
\def\labelenumi{\alph{enumi})}
\setcounter{enumi}{1}
\tightlist
\item
  Let \(q \in \mathcal{C}^2([0,1])\) satisfying \(q \geqslant 0\) . We denote \(\widetilde{\mathrm{SL}}_q\) the partial differential operator on \(\mathrm{L}^2([0,1],\mu)\) with domain
\end{enumerate}

\[
\mathrm{D} = \{f \in \mathcal {C} ^ {2} ([ 0, 1 ]) \mid f (0) = f (1) = 0 \}
\]

such that

\[
\widetilde {\mathrm{SL}} _ {q} (f) = - f ^ {\prime \prime} + q f
\]

Sturm–Liouville operator. The operator \(\widetilde{SL}_{q}\) is symmetric and essentially self-adjoint. We denote \(SL_{q}\) the closure of \(\widetilde{SL}_{q}\) .

\begin{enumerate}
\def\labelenumi{\alph{enumi})}
\setcounter{enumi}{2}
\tightlist
\item
  Let \(\lambda \in \mathrm{Sp}(\mathrm{SL}_q)\) an eigenvalue of \(\mathrm{SL}_q\) The eigenspace associated with \(\lambda\) has dimension 1 and is generated by a real-valued function belonging to D.
\end{enumerate}

\(d\) ) Every eigenvalue \(\lambda\) of \(\mathrm{SL}_q\) satisfies \(\lambda \geqslant \pi^2\) . (First prove that \(\lambda\) is positive using question \(a\) ) with \(q_{1} = q\) and \(q_{2} = 0\) .)

\begin{enumerate}
\def\labelenumi{\alph{enumi})}
\setcounter{enumi}{4}
\tightlist
\item
  For s \textgreater{} 0, there exists a unique function \(f_{s} \in \mathcal{C}^{2}([0,1])\) which is a solution of the linear differential equation
\end{enumerate}

\[
- f _ {s} ^ {\prime \prime} + q f - s ^ {2} f = 0
\]

and satisfies \(f_{s}(0)=0\) and \(f_{s}^{\prime}(0)=s\) . Set \(f(x,s)=f_{s}(x)\) ; the map f is a function of class \(C^{2}\) onto {]}0,1{[}× \(R_{+}^{*}\) . The eigenvalues of \(SL_{q}\) are of the form \(\lambda=s^{2}\) where s satisfies \(f(1,s)=0\) .

\(f)\) For \(t \in [0,1]\) and \(s > 0\) , let \(\mathrm{N}(t,s) \in \overline{\mathbf{N}}\) the number of \(x \in [0,t]\) such that \(f(x,s) = 0\) . We have \(\mathrm{N}(t,s) \in \mathbf{N}\) and the function \(s \mapsto \mathrm{N}(t,s)\) is increasing for every \(t\) fixed. (Use question \(a\) ).)

\(g)\) Let \(t \in [0,1]\) and \(s_0 > 0\) Let \(m = \mathrm{N}(t,s_0) \in \mathbf{N}\) . Denote by

\[
0 = x _ {1} (s _ {0}) <   x _ {2} (s _ {0}) <   \dots <   x _ {m} (s _ {0})
\]

the zeros of \(f_{s_0}\) in \([0,t]\) . For \(s \geqslant 0\) let

\[
0 = x _ {1} (s) <   \dots <   x _ {m} (s)
\]

the \(m\) first zeros of \(f_{s}\) in \([0,t]\) . The functions \(s\mapsto x_{j}(s)\) are of class \(\mathbf{C}^1\) and we have

\[
\partial_ {1} (x _ {j} (s), s) x _ {j} ^ {\prime} (s) = - \partial_ {2} f (x _ {j} (s), s).
\]

\begin{enumerate}
\def\labelenumi{\alph{enumi})}
\setcounter{enumi}{7}
\tightlist
\item
  The set of eigenvalues of \(SL_{q}\) is the set of values of an increasing sequence of strictly positive real numbers
\end{enumerate}

\[
0 <   \lambda_ {0} <   \lambda_ {1} <   \dots <   \lambda_ {n} <   \dots
\]

tending toward \(+\infty\) . Moreover, the eigenspace corresponding to \(\lambda_{n}\) has dimension 1; if \(\varphi_{n}\) of which is a basis, we have \(\varphi_{n} \in \mathcal{C}^{2}([0,1])\) and the function \(\varphi_{n}\) vanishes at exactly \(n\) distinct points of \([0,1]\) . (To prove that there exists an infinite number of eigenvalues of \(\mathrm{SL}_{q}\) to use the question \(a\) ) with \(q_{1} = M - s^{2}\) , where \(M \in R\) is well chosen, to verify that \(N(1, s)\) tends to \(+\infty\) as \(s \to +\infty\) .

\begin{enumerate}
\def\labelenumi{\arabic{enumi})}
\setcounter{enumi}{20}
\tightlist
\item
  We resume the notation of the preceding exercise. We denote by \(\varphi_{n}\) the unique eigenfunction of \(\mathrm{SL}_{q}\) for the eigenvalue \(\lambda_{n}\) such that \(\varphi_{n}^{\prime}(0)=\lambda_{n}\) .
\end{enumerate}

\begin{enumerate}
\def\labelenumi{\alph{enumi})}
\tightlist
\item
  Let \(\mathrm{M} = \sup \{q(x)\}\) . For every \(n \in \mathbf{N}\) , we have
\end{enumerate}

\[
\pi^ {2} (n + 1) ^ {2} \leqslant \lambda_ {n} \leqslant \pi^ {2} (n + 1) ^ {2} + \mathrm{M},
\]

and moreover

\[
\lambda_ {n} = \pi^ {2} (n + 1) ^ {2} \Bigl (1 + \mathrm{O} \Bigl (\frac {1}{n ^ {2}} \Bigr) \Bigr)
\]

as \(n \rightarrow +\infty\) .

\begin{enumerate}
\def\labelenumi{\alph{enumi})}
\setcounter{enumi}{1}
\tightlist
\item
  Let \(s > 0\) Let \(f \in \mathcal{C}^2([0,1])\) such that
\end{enumerate}

\[
- f ^ {\prime \prime} + s ^ {2} f = - q f, \qquad f (0) = 0, \quad f ^ {\prime} (0) = s.
\]

We then have \((1-s^{-1}u_{s})f=g_{s}\) , where \(g_{s}(x)=\sin(sx)\) and \(u_{s}\) is the integral operator on \(\mathcal{C}([0,1])\) given by

\[
u _ {s} (f) (x) = \int_ {0} ^ {x} \sin (s (x - y)) q (y) f (y) d \mu (y).
\]

\begin{enumerate}
\def\labelenumi{\alph{enumi})}
\setcounter{enumi}{2}
\item
  For \(s\) sufficiently large, we have \(s \notin \operatorname{Sp}(u_s)\) .
\item
  For every fixed x, we have
\end{enumerate}

\[
\varphi_ {n} (x) = \sin ((n + 1) x) + \mathrm{O} (n ^ {- 1})
\]

as \(n \rightarrow +\infty\) .

\begin{enumerate}
\def\labelenumi{\arabic{enumi})}
\setcounter{enumi}{21}
\tightlist
\item
  We resume the notation of the preceding exercise.
\end{enumerate}

\begin{enumerate}
\def\labelenumi{\alph{enumi})}
\tightlist
\item
  There exist functions \(f_{1}\) and \(f_{2}\) in \(\mathcal{C}^{2}([a,b])\) such that \(-f_{i}^{\prime\prime} + qf_{i} = 0\) and
\end{enumerate}

\[
\begin{array}{c c} f _ {1} (0) = 0, & f _ {1} ^ {\prime} (0) = 1 \\ f _ {2} (1) = 0, & f _ {2} ^ {\prime} (1) = - 1. \end{array}
\]

These functions are the unique solutions of these equations, and the function \(w = f_{1}f_{2}' - f_{1}'f_{2}\) is constant.

\begin{enumerate}
\def\labelenumi{\alph{enumi})}
\setcounter{enumi}{1}
\tightlist
\item
  Let \(g \in \mathcal{C}([0,1])\) . There exists a unique function \(\varphi \in \mathcal{C}^2([0,1])\) such that
\end{enumerate}

\[
- \varphi^ {\prime \prime} + q \varphi = g,
\]

and we have \(\varphi = c_{1}f_{1} + c_{2}f_{2}\) , where

\[
c _ {1} (x) = - \frac {1}{w} \int_ {x} ^ {1} g (y) f _ {2} (y) d \mu (y), \qquad c _ {2} (x) = - \frac {1}{w} \int_ {0} ^ {x} g (y) f _ {1} (y) d \mu (y)
\]

for all \(x \in [0,1]\) .

\begin{enumerate}
\def\labelenumi{\alph{enumi})}
\setcounter{enumi}{2}
\tightlist
\item
  The partial operator \(\mathrm{SL}_q\) is surjective; its inverse is the map \(g\mapsto \varphi\) and there exists a function \(\mathrm{G}\in \mathcal{C}([0,1]^2)\) such that
\end{enumerate}

\[
\mathrm{SL} _ {q} ^ {- 1} (g) (x) = \int_ {0} ^ {1} \mathrm{G} (x, y) g (y) d \mu (y)
\]

for all \(g \in \mathcal{C}([0,1])\) and almost all \(x \in [0,1]\) . We have \(G(x,y) = G(y,x)\) for all \((x,y) \in [0,1]^2\) .

\(d)\) The operator \(\mathrm{SL}_q^{-1}\) is compact and injective on \(\mathrm{L}^2 ([0,1],\mu)\) .

\begin{enumerate}
\def\labelenumi{\alph{enumi})}
\setcounter{enumi}{4}
\tightlist
\item
  The functions \((\varphi_{n})\) form an orthonormal basis of \(\mathrm{L}^{2}([0,1])\) , and the spectrum of \(SL_{q}\) coincides with the set of its eigenvalues.
\end{enumerate}

\begin{enumerate}
\def\labelenumi{\arabic{enumi})}
\setcounter{enumi}{22}
\tightlist
\item
  Let \(E = \ell_{\mathbf{C}}^{2}(\mathbf{N})\) and let u be the endomorphism of E defined by
\end{enumerate}

\[
u ((x _ {i}) _ {i \in \mathbf {N}}) = (0, x _ {1}, \dots , x _ {n}, \dots)
\]

for all \((x_{i})\in \mathrm{E}\)

\begin{enumerate}
\def\labelenumi{\alph{enumi})}
\item
  The spectrum of \(u\) is the closed unit disk in \(\mathbf{C}\) .
\item
  The residual spectrum of u is the open unit disk in C.
\item
  The essential spectrum of \(u\) is the unit circle in \(\mathbf{C}\) .
\end{enumerate}

\begin{enumerate}
\def\labelenumi{\arabic{enumi})}
\setcounter{enumi}{23}
\tightlist
\item
  Let E be a complex Hilbert space and u a partial symmetric operator on E. Let \(F = E \oplus E\) and let v be the partial operator on F with domain \(\text{dom}(u^{*}) \oplus \text{dom}(\overline{u})\) such that \(v(x, y) = (\overline{u}(y), u^{*}(x))\) .
\end{enumerate}

\begin{enumerate}
\def\labelenumi{\alph{enumi})}
\item
  The partial operator v is self-adjoint.
\item
  Let \(\widetilde{\mathrm{E}}\) the image of the linear map \(\delta\colon x\mapsto(x,x)\) The reduction of \(v\) to \(\widetilde{\mathrm{E}}\) is the partial operator with domain \(\delta(\mathrm{dom}(u))\) such that \((x,x)\mapsto(u(x),u(x))\) .
\end{enumerate}

